\chapter{مقدمه}
\label{ch:intro}

\section{زمینه}
\label{sec:intro-background}

نرم‌افزارهای امروزی به‌شکل فزاینده‌ای روی زیرساخت‌های رایانش ابری\LTRfootnote{Cloud Computing}
مستقر می‌شوند و با معماری ریزخدمت\LTRfootnote{Microservice} توسعه می‌یابند. در این معماری،
یک سامانه‌ی بزرگ به مجموعه‌ای از سرویس‌های کوچک و مستقل شکسته می‌شود که هر یک جداگانه مستقر و
جداگانه مقیاس‌گذاری می‌شود. همین استقلال، امکانی را فراهم می‌کند که در معماری یکپارچه وجود
نداشت: ظرفیت را می‌توان دقیقاً به همان بخشی افزود که زیر فشار است، نه به کل سامانه.

مهم‌ترین مزیتی که رایانش ابری به این معماری اضافه می‌کند، کشسانی\LTRfootnote{Elasticity}
است؛ یعنی توانایی سامانه در افزودن یا کاستن خودکار منابع پردازشی متناسب با بار ورودی. هربست و
همکاران~\cite{herbst2013} در تعریف دقیقی که برای این مفهوم پیشنهاد کرده‌اند، کشسانی را از
مقیاس‌پذیری\LTRfootnote{Scalability} و کارایی\LTRfootnote{Efficiency} تفکیک می‌کنند:
مقیاس‌پذیری یک ویژگی ایستا است و می‌گوید سامانه \emph{می‌تواند} با منابع بیشتر بار بیشتری را
سرویس دهد، اما کشسانی یک ویژگی پویاست و می‌پرسد سامانه \emph{چقدر سریع و چقدر دقیق} خود را با
تغییر بار تطبیق می‌دهد. این تمایز برای کل این نوشتار پایه‌ای است، چون آنچه در فصل‌های بعد
اندازه‌گیری می‌شود سرعت و دقت تطبیق است و نه صرفاً امکان آن.

مدیریت درست منابع در این بستر، مسئله‌ای صرفاً فنی نیست. کم‌تأمین منابع مستقیماً به افت کیفیت
خدمت\LTRfootnote{Quality of Service (QoS)} و نقض توافق سطح خدمت\LTRfootnote{Service Level
Agreement (\lr{SLA})} تبدیل می‌شود، و بیش‌تأمین منابع مستقیماً هزینه‌ی زیرساخت است. ورما و
همکاران~\cite{verma2015} در گزارش سامانه‌ی \lr{Borg} در گوگل نشان می‌دهند که در مقیاس بزرگ،
اختلاف میان منابع درخواست‌شده و منابع واقعاً مصرف‌شده به هدر رفتن حجم عظیمی از ظرفیت می‌انجامد؛
و رزادسا و همکاران~\cite{rzadca2020} در معرفی \lr{Autopilot} همین مسئله را انگیزه‌ی اصلی
خودکارسازی تعیین منابع می‌دانند. به همین دلیل، مقیاس‌گذاری خودکار\LTRfootnote{Autoscaling} به
یکی از موضوعات اصلی نگهداشت سامانه‌های ابری تبدیل شده است.

\section{مقیاس‌گذاری خودکار و انواع آن}
\label{sec:intro-autoscaling}

مقیاس‌گذاری خودکار سه شکل رایج دارد. در \emph{مقیاس‌گذاری افقی}\LTRfootnote{Horizontal Scaling}
تعداد نمونه‌های سرویس تغییر می‌کند. در \emph{مقیاس‌گذاری عمودی}\LTRfootnote{Vertical Scaling}
منابع اختصاص‌یافته به هر نمونه (پردازنده و حافظه) تغییر می‌کند. در \emph{مقیاس‌گذاری خوشه}
تعداد گره‌های زیرساخت تغییر می‌کند. این پروژه تنها به شکل نخست می‌پردازد، که رایج‌ترین شکل در
سکوهای ارکستراسیون کانتینر است و در \lr{Kubernetes} با ابزار
\lr{Horizontal Pod Autoscaler (HPA)} پیاده‌سازی شده است.

از منظر سازوکار تصمیم‌گیری، دسته‌بندی مهم‌تری وجود دارد که کل این نوشتار بر پایه‌ی آن ساخته
شده است. لوریدو-بوتران و همکاران~\cite{lorido2014} در مرور جامع خود، روش‌های مقیاس‌گذاری
خودکار را به پنج خانواده تقسیم می‌کنند: قواعد آستانه‌ای ایستا، نظریه‌ی کنترل، یادگیری تقویتی،
نظریه‌ی صف، و تحلیل سری‌های زمانی. کیو و همکاران~\cite{qu2018} نیز در مرور دیگری همین فضا را
از منظر معماری و چالش‌ها رده‌بندی می‌کنند. برای هدف این پروژه، تقسیم‌بندی ساده‌تری کافی است:

\begin{itemize}
  \item \textbf{واکنشی}\LTRfootnote{Reactive}: تصمیم بر پایه‌ی مقدار \emph{جاری} سنجه گرفته
  می‌شود. سامانه تنها \emph{پس از} تغییر بار عمل می‌کند.
  \item \textbf{پیش‌بینانه}\LTRfootnote{Predictive / Proactive}: تصمیم بر پایه‌ی مقدار
  \emph{پیش‌بینی‌شده} برای آینده‌ی نزدیک گرفته می‌شود، تا ظرفیت پیش از رسیدن بار آماده باشد.
\end{itemize}

منطق پشت رویکرد دوم روشن است و پشتوانه‌ی تجربی دارد: افزودن یک نمونه‌ی تازه آنی نیست. مائو و
هامفری~\cite{mao2012} در سنجش زمان راه‌اندازی ماشین‌های مجازی روی سه ارائه‌دهنده‌ی بزرگ ابری
نشان دادند که این تأخیر به‌قدری بزرگ است که تصمیم‌گیری واکنشی را ذاتاً عقب‌مانده می‌کند. در
دنیای کانتینرها این تأخیر کوچک‌تر است اما صفر نیست؛ در بستر آزمایش همین پروژه، یک نمونه‌ی
تازه حدود ۳۰ ثانیه طول می‌کشد تا آماده‌ی سرویس شود. اگر تصمیم مقیاس‌گذاری تنها پس از بالا رفتن
بار گرفته شود، سرویس در تمام این مدت کم‌تأمین است.

\section{بیان مسئله}
\label{sec:intro-problem}

ابزار پیش‌فرض \lr{Kubernetes} برای مقیاس‌گذاری افقی، یعنی \lr{HPA}، تعداد نمونه‌های یک سرویس
را بر پایه‌ی سنجه‌های مشاهده‌شده تنظیم می‌کند و در عمل کارآمد است. با این حال دو محدودیت دارد
که انگیزه‌ی این پروژه از آن‌ها می‌آید.

\textbf{محدودیت نخست، جنس سیگنال است.} پیکربندی متعارف \lr{HPA} بر مصرف پردازنده استوار است.
مصرف پردازنده یک سنجه‌ی \emph{زیرساختی} است و تقریبی غیرمستقیم از بار واقعی برنامه به شمار
می‌رود. برای سرویسی که کاملاً محدود به پردازنده است این تقریب خوب کار می‌کند، اما برای سرویسی
که منتظر پایگاه داده یا یک سرویس بیرونی می‌ماند، مصرف پردازنده زیر بار فزاینده تقریباً ثابت
می‌ماند و سیگنال، بار را دنبال نمی‌کند. آنچه چنین سرویسی واقعاً باید بر پایه‌ی آن مقیاس بگیرد،
سنجه‌های سطح‌برنامه\LTRfootnote{Application-level Metrics} است: نرخ درخواست بر ثانیه، تأخیر
پاسخ، یا طول صف. نگوین و همکاران~\cite{nguyen2020} در بررسی تجربی خود از \lr{HPA} دقیقاً همین
تفاوت را می‌سنجند و نشان می‌دهند که انتخاب میان سنجه‌های منبعی و سنجه‌های سفارشی
\lr{Prometheus} رفتار مقیاس‌گذار را به‌طور معناداری تغییر می‌دهد.

\lr{Prometheus} به‌عنوان استاندارد عملی جمع‌آوری سنجه در بوم‌سازگان ابری، این داده‌ها را در
اختیار می‌گذارد؛ اما به‌تنهایی سیاست مقیاس‌گذاری اعمال نمی‌کند. پل زدن میان این دو، از راه
\lr{Prometheus Adapter} و واسط سنجه‌های سفارشی \lr{Kubernetes} ممکن است، ولی همچنان سیاستِ
تصمیم‌گیری همان سیاست \lr{HPA} می‌ماند.

\textbf{محدودیت دوم، جنس سیاست است.} \lr{HPA} کاملاً واکنشی است. این محدودیت ذاتی طراحی آن
است و نه نقص پیاده‌سازی. با توجه به تأخیر راه‌اندازی نمونه‌ها که در
بخش~\ref{sec:intro-autoscaling} گفته شد، یک کنترل‌گر واکنشی همیشه به اندازه‌ی مجموع
«فاصله‌ی مشاهده + فاصله‌ی تصمیم + زمان آماده‌شدن نمونه» از بار عقب است.

راه‌حل متعارف برای محدودیت دوم، افزودن پیش‌بینی به حلقه‌ی تصمیم است. روی و
همکاران~\cite{roy2011} یکی از نمونه‌های شناخته‌شده‌ی این رویکردند و توکا و
همکاران~\cite{toka2021} آن را روی خودِ \lr{Kubernetes} به کار برده‌اند. در صنعت نیز
\lr{Netflix} با سامانه‌ی \lr{Scryer}~\cite{scryer} و ارائه‌دهندگان بزرگ ابری با قابلیت
مقیاس‌گذاری پیش‌بینانه، همین ایده را عملیاتی کرده‌اند.

\textbf{اما پرسشی که کمتر پرسیده شده، این است که افزودن پیش‌بینی به یک حلقه‌ی کنترل بسته چه
اثری بر پایداری آن حلقه دارد.} ادبیات مقیاس‌گذاری پیش‌بینانه عمدتاً کیفیت پیش‌بینی‌کننده را
به‌صورت \emph{مستقل} می‌سنجد؛ یعنی خطای پیش‌بینی را روی یک سری زمانی ثبت‌شده اندازه می‌گیرد.
این سنجش، صرف‌نظر از دقتش، یک چیز را نمی‌بیند: در سامانه‌ی واقعی، پیش‌بینی‌کننده درون حلقه
قرار دارد و ورودی‌اش را کنش‌های همان حلقه شکل می‌دهد. همان‌طور که هلراشتاین و
همکاران~\cite{hellerstein2004} در کتاب مرجع کنترل پس‌خوردی سامانه‌های محاسباتی تأکید
می‌کنند، پایداری یک حلقه‌ی بسته را نمی‌توان از روی رفتار مؤلفه‌های آن به‌تنهایی نتیجه گرفت.

\section{انگیزه و پرسش پژوهش}
\label{sec:intro-question}

کار این پروژه با ساختن یک مقیاس‌گذار خودکار سفارشی مبتنی بر \lr{Prometheus} آغاز شد. در
نخستین هفته‌ی پیاده‌سازی، اندازه‌گیری‌ای انجام شد که مسیر کار را تغییر داد و پرسش را
تیزتر کرد. جزئیات کامل آن در بخش~\ref{sec:design-perreplica} و فصل~\ref{ch:results} می‌آید،
اما خلاصه‌اش چنین است.

نمونه‌ی اولیه، بار \emph{هر نمونه} را پیش‌بینی می‌کرد. اما بار هر نمونه برابر است با بار کل
تقسیم بر تعداد نمونه‌های آماده، و تعداد نمونه‌های آماده همان کمیتی است که مقیاس‌گذار خودش
تعیین می‌کند. پس پیش‌بینی‌کننده سیگنالی را برون‌یابی می‌کرد که کنش‌های خودش آن را جابه‌جا
می‌کرد. نتیجه، یک حلقه با بهره‌ی مثبت است: مقیاس به بالا باعث افت بار هر نمونه می‌شود،
پیش‌بینی‌کننده این افت را ادامه‌دار می‌بیند، مقیاس به پایین توصیه می‌کند، و چرخه از نو آغاز
می‌شود.

این یافته دو ویژگی داشت که آن را به هسته‌ی پروژه تبدیل کرد. نخست آنکه با \emph{اندازه‌گیری}
به دست آمد و نه با مطالعه‌ی کد؛ رفتار هر مؤلفه به‌تنهایی درست بود. دوم آنکه در سنجه‌های
متعارف ارزیابی تقریباً \emph{نامرئی} بود، چون سازوکار پایدارسازی\LTRfootnote{Stabilization
Window} \mdash{} که به دلایلی کاملاً متفاوت در طراحی گنجانده شده بود \mdash{} آن را می‌پوشاند.

بر همین پایه، پرسش این پروژه از «چگونه یک مقیاس‌گذار خودکار بسازیم» به پرسش زیر تغییر یافت:

\begin{quote}
\textbf{مقیاس‌گذاری پیش‌بینانه در چه شرایطی ناپایدار می‌شود، و طراحی آن باید چگونه باشد که
ناپایدار نشود؟}
\end{quote}

این بازتعریف، سه پرسش فرعی را به دنبال دارد که ساختار ارزیابی پروژه بر آن‌ها استوار است:

\begin{enumerate}
  \item \textbf{پرسش ۱.} پیش‌بینیِ کدام سیگنال، حلقه‌ی کنترل را ناپایدار می‌کند و چرا؟ آیا
  می‌توان شرط پایداری را به شکلی ساختاری بیان کرد و نه صرفاً تجربی؟
  \item \textbf{پرسش ۲.} سازوکار پایدارسازی چه مقدار از این ناپایداری را می‌پوشاند، و پوشاندن
  آن چه بهایی دارد؟ آیا سیاستی که با میراگر روشن از همه‌ی آزمون‌ها سربلند بیرون می‌آید،
  واقعاً پایدار است؟
  \item \textbf{پرسش ۳.} در برابر مقیاس‌گذار پیش‌فرض \lr{Kubernetes}، پیش‌بینی چه چیزی
  می‌خرد و بهای آن چیست؟ این سود روی کدام جنس بار کاری وجود دارد و روی کدام جنس ذاتاً وجود
  ندارد؟
\end{enumerate}

توجه به این نکته اهمیت دارد که ساختن یک مقیاس‌گذار مبتنی بر \lr{Prometheus} به‌خودی‌خود کار
تازه‌ای نیست: \lr{KEDA} یک پروژه‌ی فارغ‌التحصیل بنیاد \lr{CNCF} است~\cite{kedagraduation}،
و ترکیب \lr{Prometheus Adapter} با \lr{HPA} همین کار را انجام می‌دهد. بنابراین در این
نوشتار، مقیاس‌گذار پیاده‌سازی‌شده \emph{ابزار} پژوهش است و نه \emph{دستاورد} آن. دستاورد،
تحلیل پایداری است.

\section{اهداف و نیازمندی‌ها}
\label{sec:intro-requirements}

ابزاری که این پژوهش به آن نیاز دارد باید نیازمندی‌های زیر را برآورده کند. این فهرست، مرز
کارکردی مقیاس‌گذارِ ساخته‌شده را تعیین می‌کند و در فصل‌های سوم تا پنجم به هر یک پاسخ داده
می‌شود. جدول~\ref{tab:intro-req-map} نشان می‌دهد هر نیازمندی کجا برآورده و کجا ارزیابی شده
است.

\subsection*{نیازمندی‌های کارکردی}

\begin{enumerate}
  \item جمع‌آوری سنجه‌ها از \lr{Prometheus} با استفاده از پرس‌وجوی \lr{PromQL}.
  \item اعمال سیاست‌های مقیاس‌گذاری بر پایه‌ی سنجه‌های سطح‌برنامه (نرخ درخواست، تأخیر، طول صف).
  \item تغییر تعداد نمونه‌های سرویس هدف از راه واسط برنامه‌نویسی \lr{Kubernetes}.
  \item پشتیبانی از سیاست‌های چندگانه و پیکربندی‌پذیر (آستانه‌ای و پیش‌بینانه) بدون تغییر در کد.
  \item انتشار سنجه‌های داخلی خودِ مقیاس‌گذار برای مشاهده‌پذیری و مصورسازی در \lr{Grafana}.
\end{enumerate}

\subsection*{نیازمندی‌های غیرکارکردی}

\begin{enumerate}
  \item زمان واکنش کوتاه به تغییرات بار.
  \item پایداری تصمیم‌ها و پرهیز از نوسان پیاپی مقیاس\LTRfootnote{Anti-flapping}.
  \item سربار منابع پایین برای خود مقیاس‌گذار.
  \item قابلیت اطمینان و در دسترس بودن.
\end{enumerate}

\begin{table}[H]
\centering
\caption{نگاشت نیازمندی‌ها به بخش‌های این نوشتار}
\label{tab:intro-req-map}
\begin{tabular}{|r|l|l|}
\hline
\textbf{نیازمندی} & \textbf{پیاده‌سازی} & \textbf{ارزیابی} \\
\hline
جمع‌آوری با \lr{PromQL} & بخش~\ref{sec:design-collector} & بخش~\ref{sec:results-integration} \\
سنجه‌های سطح‌برنامه & بخش~\ref{sec:design-collector} & بخش~\ref{sec:method-arms} \\
تغییر تعداد نمونه‌ها & بخش~\ref{sec:design-scaler} & بخش~\ref{sec:results-integration} \\
سیاست‌های چندگانه & بخش~\ref{sec:design-policy} & فصل~\ref{ch:results} \\
سنجه‌های داخلی & بخش~\ref{sec:design-observability} & بخش~\ref{sec:results-stabilizer} \\
\hline
زمان واکنش & بخش~\ref{sec:design-loop} & بخش~\ref{sec:results-reaction} \\
پرهیز از نوسان & بخش~\ref{sec:design-stabilizer} & بخش~\ref{sec:results-stabilizer} \\
سربار پایین & بخش~\ref{sec:design-loop} & بخش~\ref{sec:results-overhead} \\
قابلیت اطمینان & بخش~\ref{sec:design-loop} & بخش~\ref{sec:results-integration} \\
\hline
\end{tabular}
\end{table}

\section{دستاوردهای پروژه}
\label{sec:intro-contributions}

دستاوردهای این پروژه به ترتیب اهمیت چنین‌اند:

\begin{enumerate}
  \item \textbf{شناسایی و توضیح یک ناپایداری حلقه‌ی بسته در مقیاس‌گذاری پیش‌بینانه.} نشان داده
  می‌شود که پیش‌بینی بار \emph{هر نمونه} \mdash{} یعنی سیگنالی که کنترل‌گر خودش آن را جابه‌جا می‌کند \mdash{}
  به پس‌خور مثبت می‌انجامد، و پیش‌بینی بار \emph{کل} \mdash{} که ورودی برون‌زاست \mdash{} چنین مشکلی ندارد.
  این تفاوت هم به‌صورت ساختاری استدلال شده و هم در سه سطح مستقل اندازه‌گیری شده است: آزمون
  واحد حلقه‌بسته، شبیه‌ساز، و خوشه‌ی واقعی.

  \item \textbf{نشان دادن اینکه سازوکار پایدارسازی، ناپایداری را می‌پوشاند و نه اصلاح.} این
  نتیجه‌ی اصلی پروژه است. یک پنجره‌ی پایدارسازی ۹۰ ثانیه‌ای، یک قانون کنترل بنیاداً معیوب را
  چنان می‌پوشاند که از همه‌ی آزمون‌های کیفیت خدمت و هزینه سربلند بیرون می‌آید. با خاموش کردن
  همین میراگر، سیاست معیوب در یک اجرای ۱۸۰۰ ثانیه‌ای به ۱۱۹ تغییر جهت در ۱۲۰ چرخه‌ی کنترل و
  ۵۰/۸ درصد نقض کیفیت خدمت می‌رسد. پیامد روش‌شناختی این یافته آن است که ارزیابی پایداری باید
  میراگر را خاموش کند.

  \item \textbf{ارزیابی تجربی مقایسه‌ای بر روی خوشه‌ی واقعی}، شامل ۴۹ اجرای معتبر با مقیاس
  زمانی حقیقی، چهار الگوی بار و نُه بازوی مقایسه، از جمله دو خط مبنای ایستا در دو سر منحنی
  هزینه/کیفیت و دو بازوی \lr{HPA} استاندارد.

  \item \textbf{ابطال یک فرض رایج درباره‌ی سود پیش‌بینی.} تصور متعارف \mdash{} که فرض
  آغازین این پروژه هم بود \mdash{} آن است که مقیاس‌گذاری پیش‌بینانه به «جهش‌های ناگهانی
  ترافیک» کمک می‌کند. اندازه‌گیری خلاف آن را نشان می‌دهد: هیچ
  پیش‌بینی‌کننده‌ای یک پله‌ی آنی را از پیش نمی‌بیند. ادعای قابل دفاع آن است که پیش‌بینی روی
  بارهای \emph{روندار و دوره‌ای} کمک می‌کند و روی پله‌های پیش‌بینی‌ناپذیر ذاتاً نمی‌تواند
  کمک کند.

  \item \textbf{یک چارچوب آزمایش تکرارپذیر با بررسی اعتبار خودکار.} چارچوب پیش و پس از هر اجرا
  چندین شرط اعتبار را وارسی می‌کند و اجراهای معیوب را پیش از ورود به تحلیل کنار می‌گذارد.
  تفکیک «اندازه‌گیری» از «تحلیل» به این معناست که هر تعریف سنجه که بعداً نادرست از آب درآید،
  بدون بازگشت به خوشه قابل اصلاح و بازامتیازدهی است.

  \item \textbf{ثبت مجموعه‌ای از دام‌های خاموش} در مسیر ساخت چنین سامانه‌ای \mdash{} از پایه‌ی نادرست
  فرمول گرفته تا فروپاشی مشاهده‌پذیری زیر بار اضافی و نقص تحویل بار در مولد بار \mdash{} که همگی
  به‌جای شکست پرسروصدا، نتایج را بی‌سروصدا خراب می‌کنند.
\end{enumerate}

\section{محدوده‌ی کار}
\label{sec:intro-scope}

برای آنکه ادعاهای این نوشتار به‌درستی خوانده شوند، لازم است مرز کار صریح باشد.

\textbf{آنچه در محدوده است:} مقیاس‌گذاری افقی یک سرویس بدون‌حالت روی یک خوشه‌ی
\lr{Kubernetes}؛ سیگنال مقیاس‌گذاری از جنس نرخ درخواست خوانده‌شده از \lr{Prometheus}؛ دو
خانواده‌ی سیاست (آستانه‌ای و پیش‌بینانه)؛ دو پیش‌بینی‌کننده‌ی هموارسازی نمایی؛ و مقایسه با
\lr{HPA} استاندارد در دو پیکربندی.

\textbf{آنچه در محدوده نیست:} مقیاس‌گذاری عمودی و مقیاس‌گذاری خوشه؛ سرویس‌های حالت‌دار؛
زنجیره‌های چندسرویسی با وابستگی متقابل؛ پیش‌بینی‌کننده‌های سنگین مانند \lr{ARIMA} و
\lr{LSTM} که به‌عنوان کار آینده در بخش~\ref{sec:conclusion-future} پیشنهاد شده‌اند؛ و
خوشه‌های چندگره با رقابت منابع میان مستأجران مختلف. محدودیت‌های اعتبار اندازه‌گیری‌ها به‌طور
مفصل در بخش~\ref{sec:results-validity} آمده است.

\section{روش کار}
\label{sec:intro-method}

روش کار این پروژه سه سطحی است و در فصل~\ref{ch:method} به تفصیل شرح داده می‌شود:

\begin{itemize}
  \item \textbf{آزمون واحد} برای منطق موتور سیاست، مستقل از هر زیرساختی. به‌لطف تفکیک
  واسط‌ها، یک حلقه‌ی کنترل \emph{بسته} به‌طور کامل درون یک آزمون واحد اجرا می‌شود و دینامیک
  سیاست بدون \lr{Prometheus} و بدون خوشه سنجیده می‌شود.
  \item \textbf{آزمون یکپارچه‌سازی} برای رفتار سامانه روی خوشه‌ی واقعی.
  \item \textbf{آزمون کارایی و مقایسه} در برابر \lr{HPA} استاندارد، تحت الگوهای بار یکسان.
\end{itemize}

افزون بر این سه سطح، یک شبیه‌ساز گسسته به زبان \lr{Python} نیز نگهداری شد که همان الگوهای بار
و همان پیکربندی را اجرا می‌کند. نقش آن، اعتبارسنجی متقابل ترتیب نتایج است؛ محدودیت‌های این
اعتبارسنجی در بخش~\ref{sec:results-simulator} گزارش شده است.

\section{ساختار این نوشتار}
\label{sec:intro-structure}

\begin{itemize}
  \item \textbf{فصل~\ref{ch:background}} مفاهیم پایه را معرفی می‌کند: \lr{Kubernetes} و
  الگوریتم دقیق \lr{HPA}، \lr{Prometheus} و \lr{PromQL}، روش‌های پیش‌بینی سری زمانی، و
  مفاهیم لازم از کنترل پس‌خوردی. سپس کارهای پیشین مرور و شکاف پژوهشی این پروژه مشخص می‌شود.

  \item \textbf{فصل~\ref{ch:design}} طراحی و پیاده‌سازی مقیاس‌گذار را شرح می‌دهد: معماری
  سه‌بخشی، جمع‌آورنده‌ی سنجه، موتور سیاست و ریاضیات پیش‌بینی‌کننده‌ها، سازوکار پایدارسازی،
  حلقه‌ی کنترل، و نقص‌هایی که نوشتن آزمون آشکار کرد.

  \item \textbf{فصل~\ref{ch:method}} روش‌شناسی ارزیابی را تعریف می‌کند: بستر آزمایش،
  کالیبراسیون برنامه‌ی نمونه، الگوهای بار، بازوهای مقایسه، تعریف دقیق سنجه‌ها، و چارچوب
  آزمایش با بررسی‌های اعتبار آن.

  \item \textbf{فصل~\ref{ch:results}} نتایج را گزارش و تحلیل می‌کند، شامل نتیجه‌ی اصلی
  پروژه درباره‌ی اثر سازوکار پایدارسازی و بخشی مفصل درباره‌ی محدودیت‌های اعتبار.

  \item \textbf{فصل~\ref{ch:conclusion}} جمع‌بندی، پاسخ به پرسش‌های پژوهش، درس‌های مهندسی و
  کارهای آینده را ارائه می‌کند.
\end{itemize}

پیوست‌ها شامل پیکربندی کامل، راهنمای بازتولید آزمایش‌ها، و جدول‌های کامل نتایج به تفکیک هر
الگوی بار است.
